\subsection*{从一个熟悉但尚未解释的公式开始}
圆面积是 $\pi R^2$，但若只背公式，就不知道它为什么与半径平方有关。把圆按距离中心的远近分成薄圆环：半径为 $r$ 的那一层，长度约为周长 $2\pi r$，厚度为一个小增量。越外层越长，因此总面积是在累积不断增加的周长。

有限厚度的圆环并不等于长方形。外半径比内半径大 $h$ 时，精确面积是 $2\pi rh+\pi h^2$；长方形近似少算了 $\pi h^2$。当我们把厚度切细，总误差能否趋零，才是从近似过渡到公式的关键。这也是极限需要出场的原因。

\subsection*{把累积的问题倒过来}
现在假设面积函数已经知道，半径稍微增加时，面积增加多快？用新增面积除以 $h$，得到 $2\pi r+\pi h$，它趋向周长。于是“把周长累积起来”和“观察面积的瞬时变化”互相连接，这就是积分与导数关系的原型。

先在更简单的直线高度 $f(x)=x$ 上练习。到 $x$ 为止的三角形面积是 $x^2/2$，多出宽度 $h$ 时，新增矩形与小三角形的面积为 $xh+h^2/2$。除以宽度后，平均高度趋向边界高度 $x$。

看面积实验时，分别注意矩形总和与真实曲线间的差，不只看矩形数量增加。问题不是“小块很多所以看起来差不多”，而是能否让误差随切分变细而受控。后面的章节将先解释导数，再回到这一累积与变化的连接。
